% 6-glosar.tex starts here
% Necesită în preambul: \usepackage{hyphenat}
% ----------------------------------------------------------
% ACRONYMS (used with \gls/\acr... rendered with \printglossary[type=\acronymtype])
% ----------------------------------------------------------

\newacronym{pan}{PAN}{Prisma–Analog–Nx}

\newacronym{ssr}{SSR}{randare pe server}
\newacronym{ssg}{SSG}{generare statică de pagini}
\newacronym{csr}{CSR}{randare pe client}

\newacronym{trpc}{tRPC}{apeluri de proceduri la distanță tipizate pentru TypeScript}

\newacronym[longplural={obiecte de transfer de date}, shortplural={DTO-uri}]
  {dto}{DTO}{obiect de transfer de date}

\newacronym[longplural={mapări obiect–relaționale}, shortplural={ORM-uri}]
  {orm}{ORM}{mapare obiect–relațională}


\newacronym{e2e}{E2E}{testare cap-coadă}

\newacronym{ci}{CI}{integrare continuă}
\newacronym{cicd}{CI/CD}{integrare continuă și livrare/implementare continuă}

\newacronym{p50}{p50}{percentila 50 (mediană)}
\newacronym{p95}{p95}{percentila 95}

\newacronym[longplural={produse minime viabile}, shortplural={MVP-uri}]
  {mvp}{MVP}{produs minim viabil}

\newacronym{sql}{SQL}{limbaj de interogare structurat}

\newacronym[longplural={interfețe de programare a aplicațiilor}, shortplural={API-uri}]
  {api}{API}{interfață de programare a aplicațiilor}

\newacronym[longplural={interfețe de utilizator}, shortplural={UI-uri}]
  {ui}{UI}{interfață de utilizator}

\newacronym[longplural={baze de date}, shortplural={DB-uri}]
  {db}{DB}{bază de date}


\newacronym[longplural={localizatori uniformi de resurse}, shortplural={URL-uri}]
  {url}{URL}{localizator uniform de resurse}

\newacronym{mean}{MEAN}{MongoDB–Express–Angular–Node.js}


\newacronym{html}{HTML}{HyperText Markup Language (limbaj de marcare hipertext)}

\newacronym{spa}{SPA}{aplicație web de o singură pagină (single-page application)}
\newacronym{seo}{SEO}{optimizare pentru motoarele de căutare}
\newacronym[longplural={cereri de tragere}, shortplural={PR-uri}]
  {pr}{PR}{cerere de integrare (pull request)}


\newacronym{http}{HTTP}{protocol de transfer hipertext}


\newacronym[longplural={modele ale obiectelor documentului}, shortplural={DOM-uri}]
  {dom}{DOM}{model al obiectelor documentului}


% ----------------------------------------------------------
% TERMS and NAMES (used with \gls rendered with \printglossary[title={Glosar de termeni}])
% ----------------------------------------------------------

\newglossaryentry{web}{
  name={\NoHyphens{web}},
  sort={web},
  description={ecosistemul de resurse \emph{hipertext} accesibile prin Internet (World Wide Web); în lucrare, „aplicație web” desemnează aplicație accesată din browser, cu randare \glsxtrshort{ssr}/\glsxtrshort{ssg}/\glsxtrshort{csr}},
}

\newglossaryentry{nx}{
  name={\NoHyphens{Nx}},
  description={sistem de build și orchestrare pentru monorepo-uri JavaScript/TypeScript},
}

\newglossaryentry{prisma}{
  name={\NoHyphens{Prisma}},
  description={bibliotecă \glsxtrshort{orm} pentru TypeScript/Node.js},
}

\newglossaryentry{analog}{
  name={\NoHyphens{Analog}},
  description={meta-framework full-stack pentru \gls{angular}, cu suport \glsxtrshort{ssr} și rutare pe fișiere; folosit pentru aplicații \gls{web}},
}

\newglossaryentry{angular}{
  name={\NoHyphens{Angular}},
  description={platformă și framework pentru aplicații \gls{web} (SPA/\glsxtrshort{ssr})},
}

\newglossaryentry{nitro}{
  name={\NoHyphens{Nitro}},
  description={runtime universal (UnJS) pentru rute server și randare pe server (\glsxtrshort{ssr})},
}


\newglossaryentry{zod}{
  name={\NoHyphens{Zod}},
  description={bibliotecă de validare pe bază de schemă pentru TypeScript/JavaScript},
}

\newglossaryentry{cockroachdb}{
  name={\NoHyphens{CockroachDB}},
  description={sistem de gestiune a bazelor de date distribuit (NewSQL), compatibil la nivel de protocol cu \gls{postgresql}, cu replicare și partiționare automate pentru reziliență și scalare orizontală},
  sort={cockroachdb},
}

\newglossaryentry{postgresql}{
  name={\NoHyphens{PostgreSQL}},
  description={sistem de gestiune a bazelor de date relaționale open-source, cu suport extins pentru standardul \glsxtrshort{sql}, tranzacții ACID și extensibilitate prin tipuri și funcții definite de utilizator},
  sort={postgresql},
}


\newglossaryentry{avg}{
  name={\NoHyphens{avg}},
  sort={avg},
  description={medie aritmetică a unui set de valori; în SQL corespunde funcției \texttt{AVG}, care (în implementările uzuale) ignoră valorile \texttt{NULL}. Se folosește în raportare împreună cu percentile precum \gls{p50} (mediană) și \gls{p95}},
}

\newglossaryentry{vercel}{
  name={\NoHyphens{Vercel}},
  description={platformă cloud pentru build, preview și deploy de aplicații \gls{web}, cu integrare \glsxtrshort{cicd}, variabile de mediu per mediu și suport pentru \glsxtrshort{ssr}, funcții \gls{serverless} și \gls{edgefunctions}},
  sort={vercel},
}

\newglossaryentry{vitest}{
  name={\NoHyphens{Vitest}},
  description={framework de testare rapid pentru ecosistemul Vite, orientat spre teste unitare și de integrare, cu API în mare parte compatibil Jest și suport nativ pentru TypeScript},
  sort={vitest},
}

\newglossaryentry{playwright}{
  name={\NoHyphens{Playwright}},
  description={bibliotecă pentru automatizarea browserelor și testare \glsxtrshort{e2e}, cu suport pentru Chromium, Firefox și WebKit, rulare headless/GUI și unelte precum trace viewer},
  sort={playwright},
}

\newglossaryentry{edgefunctions}{
  name={\NoHyphens{Edge Functions}},
  description={funcții executate în rețeaua de \emph{edge} (aproape de utilizator), utile pentru latență redusă, personalizare și pre-/post-procesare a răspunsurilor \glsxtrshort{ssr}; pot rula cu constrângeri mai stricte (timp CPU, memorie) față de funcțiile serverless clasice},
  sort={edgefunctions},
}

\newglossaryentry{serverless}{
  name={\NoHyphens{serverless}},
  description={model operațional în care alocarea resurselor și scalarea sunt gestionate de platformă; dezvoltatorul furnizează funcții sau servicii ce pornesc la cerere, cu facturare pe utilizare; folosit adesea pentru \glsxtrshortpl{api}, \glsxtrshort{ssr} și joburi programate},
  sort={serverless},
}

\newglossaryentry{hydration}{
  name={\NoHyphens{hidratare (hydration)}},
  description={procesul prin care codul din client atașează comportament interactiv la \emph{markup}-ul generat prin \glsxtrshort{ssr}; permite trecerea de la randare inițială pe server la interacțiuni \glsxtrshort{csr} fără reîncărcarea paginii},
  sort={hydration},
}

\newglossaryentry{monorepo}{
  name={\NoHyphens{monorepo}},
  description={strategie de organizare a codului în care mai multe aplicații și biblioteci coexistă într-un singur depozit, cu orchestrare unitară a build-urilor, testelor și \glsxtrshort{cicd}; în contextul lucrării este gestionat cu Nx},
  sort={monorepo},
}


\newglossaryentry{typescript}{
  name={\NoHyphens{TypeScript}},
  sort={typescript},
  description={superset tipizat al limbajului JavaScript care introduce tipuri statice și se transpilează la JavaScript},
}

\newglossaryentry{kookio}{
  name={\NoHyphens{Kookio}},
  sort={kookio},
  description={aplicație \gls{web} full-stack folosită ca studiu de caz în această lucrare pentru gestiunea cămării, planificarea zilnică a meselor și derivarea automată a listei de cumpărături, construită pe arhitectura \gls{pan}},
}

\newglossaryentry{lualatex}{
  name={\NoHyphens{LuaLaTeX}},
  sort={lualatex},
  description={motor \LaTeX{} bazat pe LuaTeX, cu suport nativ pentru Unicode și fonturi OpenType},
}

\newglossaryentry{biblatex}{
  name={\NoHyphens{biblatex}},
  sort={biblatex},
  description={pachet \LaTeX{} pentru gestionarea bibliografiilor și citărilor cu backend-uri moderne (de ex. \texttt{biber})},
}

\newglossaryentry{unicode}{
  name={\NoHyphens{Unicode}},
  sort={unicode},
  description={standard de codificare a caracterelor care oferă o mapare unică pentru simboluri folosite în scrieri și tehnică (ex. caractere \emph{box-drawing})},
}


\newglossaryentry{nextjs}{
  name={\NoHyphens{Next.js}},
  sort={nextjs},
  description={framework React orientat pe randare hibridă (\glsxtrshort{ssr}/\glsxtrshort{ssg}), frecvent asociat cu JAMstack},
}


\newglossaryentry{freemono}{
  name={\NoHyphens{FreeMono}},
  sort={freemono},
  description={font monospațiat din familia GNU FreeFont, cu suport Unicode extins (inclusiv caractere \emph{box-drawing} U+2500–U+257F)},
}

\newglossaryentry{treeverb}{
  name={\NoHyphens{TreeVerb}},
  sort={treeverb},
  description={mediu \emph{verbatim} personalizat (bazat pe \texttt{fvextra}/\texttt{fancyvrb}) pentru randarea arborilor/structurilor de fișiere; folosește \gls{freemono} pentru simboluri \emph{box-drawing}},
}

\newglossaryentry{vite}{
  name={\NoHyphens{Vite}},
  sort={vite},
  description={instrument de build și server de dezvoltare pentru aplicații \gls{web} moderne, bazat pe module ECMAScript, cu pornire foarte rapidă și \emph{hot module replacement} (HMR); stă la baza \gls{vitest} și este folosit de \gls{analog} pentru \glsxtrshort{ssr} și build},
}

\newglossaryentry{pino}{
  name={\NoHyphens{Pino}},
  sort={pino},
  description={bibliotecă de logare pentru Node.js, orientată pe performanță (JSON logs, output low-overhead)},
}


\newglossaryentry{repository}{
  name={\NoHyphens{repository}},
  sort={repository},
  plural={\NoHyphens{repository-uri}},
  description={spațiu logic pentru stocarea codului sursă, a istoricului de versiuni și a metadatelor; poate fi local sau la distanță (remote, ex.: GitHub/GitLab). În această lucrare, monorepo-ul reunește mai multe proiecte într-un singur repository.},
  see={[vezi și]{monorepo}},
}

\newglossaryentry{github}{
  name={\NoHyphens{GitHub}},
  sort={github},
  description={platformă de găzduire și colaborare pentru proiecte software bazate pe sistemul de versiuni Git; oferă găzduire de \glspl{repository}, mecanisme de colaborare (pull requests, code review), management de \emph{issues} și \glsxtrshort{cicd} prin GitHub Actions; utilizată frecvent pentru automatizări și publicare.},
  see={[vezi și]{repository,cicd}},
}

\newglossaryentry{gitlab}{
  name={\NoHyphens{GitLab}},
  sort={gitlab},
  description={platformă de găzduire și colaborare pentru proiecte software bazate pe Git, disponibilă ca serviciu cloud sau instalare self-hosted; integrează nativ \glsxtrshort{cicd} (GitLab CI/CD), management de \emph{issues}, \emph{merge requests} și registru de containere; alternativă la \gls{github} cu accent pe lanțul DevOps unificat.},
  see={[vezi și]{repository,cicd}},
}

\newglossaryentry{workflow}{
  name={\NoHyphens{workflow}},
  sort={workflow},
  plural={\NoHyphens{workflow-uri}},
  description={configurație declarativă care descrie o succesiune de pași automatizați (build, testare, lint, deploy) declanșați de evenimente (ex.: push, pull request, programare). În ecosisteme precum \gls{github} Actions și \gls{gitlab} \glsxtrshort{ci}, un workflow este compus din job-uri și pași, poate folosi matrici, cache, artefacte, secrete/variabile și reguli de concurență. În acest proiect, workflow-urile orchestrează \glsxtrshort{cicd} și aliniază execuția la țintele \gls{nx}.},
  see={[vezi și]{cicd,github,gitlab,nx}},
}

\newglossaryentry{compositeaction}{
  name={\NoHyphens{composite action}},
  sort={compositeaction},
  plural={\NoHyphens{composite actions}},
  description={tip de acțiune reutilizabilă în \gls{github} Actions care grupează mai mulți pași (\texttt{run}/\texttt{uses}) într-o singură acțiune; se distribuie ca director cu fișier \texttt{action.yml}, acceptă \textit{inputs}, \textit{outputs} și \textit{secrets} și poate fi apelată dintr-un \gls{workflow} prin \texttt{uses:}. Utilă pentru a elimina duplicarea (DRY) în pipeline-urile \glsxtrshort{cicd}; limitări uzuale: nu rulează în containere proprii și nu pot defini \textit{services}.},
  see={[vezi și]{workflow,github,cicd}},
}


\newglossaryentry{redis}{
  name={\NoHyphens{Redis}},
  sort={redis},
  description={magazin de date \emph{in-memory} tip cheie–valoare, cu structuri avansate (liste, seturi, hash-uri, mulțimi ordonate, \textit{streams}) și latență foarte mică; folosit frecvent pentru caching, stocare sesiuni, cozi/mesagerie (Pub/Sub, Streams), rate-limiting și \textit{leaderboards}. Suportă persistență opțională (RDB/AOF), replicare, \textit{cluster} pentru partiționare și, de regulă, completează o bază \gls{sql}/\gls{db} principală în arhitecturi microservicii și \gls{serverless}/\gls{edgefunctions}.},
  see={[vezi și]{sql,db,serverless,edgefunctions}},
}

\newglossaryentry{framework}{
  name={\NoHyphens{framework}},
  sort={framework},
  description={set de convenții, componente și unelte care oferă o structură impusă pentru dezvoltarea de aplicații. Un \textit{framework} controlează fluxul principal al aplicației (inversiune a controlului), definește puncte de extensie și impune un mod coerent de organizare a codului; spre deosebire de o bibliotecă, pe care aplicația o apelează explicit. Exemple în lucrare: \gls{angular}; \gls{analog} funcționează ca \emph{meta-framework} peste \gls{angular} pentru \glsxtrshort{ssr} și rute server.},
  see={[vezi și]{angular,analog,ssr}},
}

\newglossaryentry{metaframework}{
  name={\NoHyphens{meta-framework}},
  sort={metaframework},
  description={framework care extinde un framework de bază pentru capabilități full-stack (ex. rutare pe fișiere, \glsxtrshort{ssr}/\glsxtrshort{ssg}, integrare cu \glsxtrshort{trpc}). În această lucrare, \gls{analog} extinde \gls{angular} peste runtime-ul \gls{nitro}.},
  see={[vezi și]{analog,angular,nitro,ssr,ssg,trpc}},
}


\newglossaryentry{randare}{
  name={\NoHyphens{randare}},
  sort={randare},
  description={procesul prin care o aplicație generează HTML (și actualizări ale interfeței) din cod și date. În web-ul modern poate avea loc pe server (\glsxtrshort{ssr}), la build-time prin \glsxtrshort{ssg}/prerender, sau în browser (\glsxtrshort{csr}); de regulă, este urmată de \gls{hydration} pentru interactivitate. În \gls{analog}, strategia de randare poate fi aleasă per rută.},
  see={[vezi și]{ssr,ssg,csr,hydration,analog}},
}


\newglossaryentry{frontend}{
  name={\NoHyphens{frontend}},
  plural={\NoHyphens{frontend-uri}},
  sort={frontend},
  description={stratul aplicației care rulează în browser și gestionează \glsxtrshort{ui}-ul, logica de interacțiune și actualizările în \glsxtrshort{csr}. În arhitecturi hibride consumă HTML-ul generat prin \glsxtrshort{ssr}/\glsxtrshort{ssg} și trece prin \gls{hydration}. Comunică cu \gls{backend}-ul prin \glsxtrshort{api} (ex.: \glsxtrshort{trpc}). În acest proiect, frontend-ul este implementat cu \gls{angular} via \gls{analog}.},
  see={[vezi și]{ui,csr,ssr,hydration,analog,trpc}},
}

\newglossaryentry{backend}{
  name={\NoHyphens{backend}},
  plural={\NoHyphens{backend-uri}},
  sort={backend},
  description={stratul server-side care implementează logica de business, expune \glsxtrshort{api}-uri și mediază accesul la date (ex.: \gls{prisma} către \gls{cockroachdb}). Poate rula ca \glsxtrshort{ssr}, funcții \gls{serverless}/\gls{edgefunctions} sau servicii dedicate. În acest proiect, backend-ul este integrat în \gls{analog} (Nitro routes și routere \glsxtrshort{trpc}).},
  see={[vezi și]{api,ssr,serverless,prisma,cockroachdb,analog,trpc}},
}

\newglossaryentry{fullstack}{
  name={\NoHyphens{full-stack}},
  plural={\NoHyphens{full-stack}},
  sort={fullstack},
  description={abordare în care aceiași dezvoltatori/proiect acoperă cap-coadă \gls{frontend}-ul și \gls{backend}-ul, inclusiv modelarea datelor și pipeline-urile de \glsxtrshort{cicd}. Beneficii: tipuri partajate, coerență arhitecturală și viteză de livrare. În această lucrare, modelul full-stack este realizat prin arhitectura \gls{pan} într-un \gls{monorepo} Nx.},
  see={[vezi și]{frontend,backend,pan,monorepo,cicd,nx}},
}

\newglossaryentry{logging}{
  name={\NoHyphens{jurnalizare (logging)}},
  sort={logging},
  description={colectarea și înregistrarea evenimentelor aplicației (mesaje structurate pe niveluri, corelare prin identificatori, export către sisteme de agregare) pentru depanare și observabilitate},
}


\newglossaryentry{darwinarm64}{
  name={\NoHyphens{darwin-arm64}},
  sort={darwinarm64},
  description={etichetă de platformă (target) utilizată la generarea binary-uri pentru ecosistemul Node.js pe macOS (kernel \textit{Darwin}) cu arhitectură ARM64 (Apple Silicon). În \gls{prisma}, selectarea acestui target indică faptul că \emph{engines} native sunt construite pentru dezvoltare locală pe macOS ARM și nu sunt compatibile cu Linux de producție; diferențele de CPU și biblioteci de sistem pot produce erori la runtime.},
}

\newglossaryentry{rhelopenssl30x}{
  name={\NoHyphens{rhel-openssl-3.0.x}},
  sort={rhelopenssl30x},
  description={etichetă de platformă pentru distribuții Linux compatibile Red Hat Enterprise Linux care folosesc biblioteca OpenSSL 3.0.x. În \gls{prisma}, acest target este folosit în \texttt{binaryTargets} pentru a genera binary-uri compatibile cu mediile de execuție tipice de producție (ex.: containere bazate pe \textit{glibc}); alegerea corectă elimină erorile de tip \texttt{wrong platform for query engine} la deploy.},
}

\newglossaryentry{prismaclient}{
  name={\NoHyphens{Prisma Client}},
  sort={prismaclient},
  description={client \texttt{TypeScript} pentru baza de date, generat automat de \gls{prisma} din fișierul \texttt{schema.prisma}; oferă un API tipizat pentru operații \glsxtrshort{sql} (CRUD, filtre, relații, tranzacții) asupra modelelor definite. Utilizează motoare native (Query Engine) distribuite ca binary-uri dependente de platformă (ex.: \gls{darwinarm64}, \gls{rhelopenssl30x}), motiv pentru care se generează pe mediul țintă de deploy; se invocă prin \texttt{npx prisma generate}.},
  see={[vezi și]{prisma,orm,cockroachdb,postgresql,darwinarm64,rhelopenssl30x}},
}

% 6-glosar.tex ends here